\documentclass[11pt, a4paper]{article}

\usepackage[utf8]{inputenc}
\usepackage[T1]{fontenc}
\usepackage{amsmath, amssymb}
\usepackage{geometry}
\usepackage{booktabs}
\usepackage{array}
\usepackage{colortbl}
\usepackage{xcolor}
\usepackage{titlesec}
\usepackage{enumitem}
\usepackage{fancyhdr}
\usepackage{mdframed}

\geometry{top=2.2cm, bottom=2.2cm, left=2.3cm, right=2.3cm}

\definecolor{ucblue}{HTML}{1a237e}
\definecolor{ucbluelight}{HTML}{3949ab}
\definecolor{rowgray}{HTML}{e8eaf6}
\definecolor{codebg}{HTML}{f3f4f6}
\definecolor{accentred}{HTML}{c62828}
\definecolor{accentgreen}{HTML}{2e7d32}
\definecolor{accentamber}{HTML}{ef6c00}

\titleformat{\section}{\large\bfseries\color{ucblue}}{}{0em}{}[\vspace{-4pt}\rule{\linewidth}{0.8pt}]
\titlespacing{\section}{0pt}{12pt}{5pt}

\titleformat{\subsection}{\normalsize\bfseries\color{ucbluelight}}{}{0em}{}
\titlespacing{\subsection}{0pt}{8pt}{3pt}

\newmdenv[
  linecolor=ucbluelight, linewidth=1.5pt,
  topline=false, rightline=false, bottomline=false,
  leftmargin=8pt, innerleftmargin=10pt,
  innertopmargin=5pt, innerbottommargin=5pt,
  backgroundcolor=rowgray
]{nota}

\pagestyle{fancy}
\fancyhf{}
\lhead{\small\color{ucblue}\textbf{ICS1113} Optimización · Informe 3}
\rhead{\small\color{ucblue}Análisis de sensibilidad · Grupo 11}
\cfoot{\small\thepage}
\renewcommand{\headrulewidth}{0.4pt}

\setlength{\parindent}{0pt}
\setlength{\parskip}{4pt}

\begin{document}

\begin{center}
  {\LARGE\bfseries\color{ucblue} Análisis de sensibilidad del modelo}\\[3pt]
  {\normalsize Comparación de escenarios sobre donaciones y penalización por vencimiento}
\end{center}
\vspace{-2pt}
{\color{ucblue}\rule{\linewidth}{1.5pt}}

\section{Resumen ejecutivo}

Se evaluaron dos modificaciones sobre los parámetros del modelo base para entender qué tan sensible es el resultado a la calibración y a la oferta de donaciones. Las tres corridas mantienen idénticos: hospitales, demanda, inventario inicial, capacidades, $p = 1000$, $L = 42$, $T = 365$.

\begin{center}
\begin{tabular}{lccc}
\toprule
\rowcolor{rowgray}\textbf{Indicador} & \textbf{Base} & \textbf{$-$5\% donaciones} & \textbf{$\eta = 200$}\\
\rowcolor{rowgray} & $(f{=}1, \eta{=}80)$ & $(f{=}0{,}95, \eta{=}80)$ & $(f{=}1, \eta{=}200)$\\
\midrule
Valor objetivo & 518.181 & 880.855 & 1.020.742 \\
Demanda satisfecha & \textbf{\color{accentgreen}100,0\%} & 99,0\% & \textbf{\color{accentgreen}100,0\%} \\
Demanda insatisfecha & 0 & \textbf{\color{accentred}542} & 0 \\
\textbf{Vencidas total} & \textbf{4.577} & \textbf{\color{accentgreen}1.637} & 3.647 \\
\quad Vencidas en banco & 1.591 & 54 & 371 \\
\quad Vencidas en hospitales & 2.986 & 1.583 & 3.276 \\
Tasa desperdicio (\% oferta) & 7,8\% & \textbf{\color{accentgreen}3,1\%} & 6,3\% \\
Tiempo de cómputo & 35 min & 19 min & 15 min \\
\bottomrule
\end{tabular}
\end{center}

\section{Configuración de los escenarios}

\textbf{Escenario base.} Datos del CSV sin modificación: 57.818 donaciones anuales, $\eta = 80$.

\textbf{Escenario A ($-$5\% donaciones).} Se aplicó \texttt{factor\_donaciones = 0.95} a todas las series $b_{g,t}$, simulando una reducción uniforme del 5\% en la planificación de donaciones. Total: 53.389 bolsas (\(-\)4.429).

\textbf{Escenario B ($\eta = 200$).} Se reemplazó la penalización por vencimiento del CSV ($\eta = 80$) por $\eta = 200$. El ratio $\eta/p$ pasa de 8\% a 20\%, dándole más peso a evitar desperdicio.

\section{Desglose del objetivo por término}

\begin{center}
\begin{tabular}{lrrr}
\toprule
\rowcolor{rowgray}\textbf{Término} & \textbf{Base} & \textbf{$-$5\% don.} & \textbf{$\eta=200$}\\
\midrule
1. Calidad de uso $\sum[(L-u)+\tau_h+\alpha_g]\,y$ & 152.021 & 207.895 & 291.295 \\
2. Demanda insatisfecha $p\sum s$ & 0 & 542.000 & 0 \\
3. Vencimientos $\eta\sum w$ & 366.160 & 130.960 & 729.447 \\
\midrule
\textbf{Total} & \textbf{518.181} & \textbf{880.855} & \textbf{1.020.742} \\
\bottomrule
\end{tabular}
\end{center}

El término 1 sube en ambos escenarios respecto a la base: el modelo se ve forzado a usar bolsas más viejas (\(L-u\) mayor) para evitar vencimientos. En el escenario B este efecto es más fuerte porque la presión por no desperdiciar es mayor.

\section{Hallazgos por escenario}

\subsection{Escenario A: $-$5\% en donaciones}

\textbf{Vencimientos bajan 64\%} (de 4.577 a 1.637), pero \textbf{aparece insatisfacción del 1,0\%} (542 bolsas).

La insatisfacción se concentra casi exclusivamente en tipos Rh negativos:

\begin{center}
\begin{tabular}{lrrr}
\toprule
\rowcolor{rowgray}\textbf{Tipo} & \textbf{Insatisf.} & \textbf{Demanda} & \textbf{\%}\\
\midrule
O$-$ & 415 & 1.778 & 23,3\% \\
A$-$ & 81 & 755 & 10,7\% \\
B$-$ & 45 & 259 & 17,4\% \\
AB$-$ & 1 & 62 & 1,6\% \\
\midrule
Rh$+$ (todos) & 0 & 49.805 & 0\% \\
\bottomrule
\end{tabular}
\end{center}

\textbf{Causa estructural:} los Rh$-$ ya tenían donaciones diarias muy bajas. Al multiplicar por 0,95 y truncar a entero, se pierde una bolsa por día por tipo:
$$8 \times 0{,}95 = 7{,}6 \to 7 \quad (\text{O$-$ pierde 365 bolsas al año})$$

\subsection{Escenario B: $\eta = 200$}

\textbf{Vencimientos bajan 20\%} (de 4.577 a 3.647) \textbf{sin sacrificar cobertura}. El cambio más interesante está en \emph{dónde} ocurren los vencimientos:

\begin{center}
\begin{tabular}{lrr}
\toprule
\rowcolor{rowgray}\textbf{Ubicación} & \textbf{Base} & \textbf{$\eta=200$}\\
\midrule
Vencidas en banco & 1.591 & \textbf{\color{accentgreen}371\ ($-$77\%)} \\
Vencidas en hospitales & 2.986 & 3.276 ($+$10\%) \\
\bottomrule
\end{tabular}
\end{center}

Con mayor $\eta$, el modelo \textbf{redistribuye agresivamente bolsas del banco a hospitales} para que tengan más oportunidades de uso antes de vencer. El banco deja de ser cementerio de bolsas, pero los hospitales absorben un poco más del desperdicio residual.

\textbf{Por tipo, la mejora se concentra en los más restringidos:}

\begin{center}
\begin{tabular}{lrrr}
\toprule
\rowcolor{rowgray}\textbf{Tipo} & \textbf{Vencidas base} & \textbf{Vencidas $\eta=200$} & \textbf{Reducción}\\
\midrule
AB$+$ (1 receptor compatible) & 617 & 386 & \textbf{$-$37\%} \\
B$+$ (2 receptores) & 1.503 & 1.133 & $-$25\% \\
A$+$ (2 receptores) & 1.026 & 808 & $-$21\% \\
O$+$ (4 receptores) & 1.172 & 1.082 & $-$8\% \\
Rh$-$ (varios receptores) & 259 & 239 & $-$8\% \\
\bottomrule
\end{tabular}
\end{center}

A menor flexibilidad de compatibilidad, mayor el rescate al subir $\eta$. Tiene sentido: las bolsas \emph{stuck} (AB$+$) son las que más se beneficiaban del rescate forzado por la penalización.

\section{Conclusiones}

\begin{nota}
La principal hipótesis previa --- que los vencimientos eran 100\% estructurales por sobre-oferta --- resultó \textbf{parcialmente equivocada}. La calibración del CSV ($\eta = 80$) infrapondera el costo del desperdicio.
\end{nota}

\begin{enumerate}[leftmargin=*, itemsep=4pt]
  \item \textbf{Reducir donaciones uniformemente NO es buena política.} Aunque baja vencimientos drásticamente, introduce escasez en los tipos Rh$-$ que ya estaban al límite. El recorte debería ser \emph{selectivo} (solo Rh$+$).
  \item \textbf{La calibración del CSV ($\eta = 80$) es demasiado permisiva.} Subirla a 200 reduce vencimientos 20\% sin costo de cobertura. Se sugiere calibrar formalmente el ratio $\eta/p$ en consulta con expertos clínicos del banco.
  \item \textbf{Hay vencimientos efectivamente estructurales:} aún con $\eta = 200$ vencen 3.647 bolsas, principalmente por el supuesto 10 (inventario inicial todo a $u=L$) que concentra 760 vencimientos en el día 43, y por sobre-oferta crónica en tipos Rh$+$.
  \item \textbf{El pico del día 43 es invariante} a ambos cambios. Sigue venciendo $\sim$760 bolsas. Solo se eliminaría reformulando el supuesto 10.
\end{enumerate}

\section{Próximos análisis sugeridos}

Los siguientes escenarios no se corrieron aún y aportarían argumentos adicionales:

\begin{center}
\begin{tabular}{p{4.8cm}p{5.5cm}p{4.5cm}}
\toprule
\rowcolor{rowgray}\textbf{Escenario} & \textbf{Predicción} & \textbf{Valor para informe}\\
\midrule
\textbf{Recorte selectivo $-$10\% solo en Rh$+$} & Vencidas $\sim$1.800--2.200, cobertura 100\% & \textbf{Alto:} recomendación realista \\
\midrule
Combinado $f{=}0{,}95$ + $\eta{=}200$ & Vencidas $\sim$1.300--1.500, insatisf. $\sim$542 & Medio: confirma aditividad \\
\midrule
Inventario inicial escalonado en edades 1..L & Elimina pico día 43, vencidas $\sim$3.500--4.000 & Medio: defiende supuesto 10 \\
\bottomrule
\end{tabular}
\end{center}

\textbf{Recomendación:} priorizar el \emph{recorte selectivo Rh$+$}. Es la única política que combina las ventajas de A (menos desperdicio) y B (100\% cobertura) sin sus contras.

\vfill
{\color{ucblue}\rule{\linewidth}{0.6pt}}
\begin{center}
  \small Tres corridas Gurobi · misma instancia (8 hospitales, 8 tipos, $T=365$, $L=42$)\\
  \small Modelo PDF v2 (R1..R15) · óptimo global en cada escenario
\end{center}

\end{document}
